% !TeX program = pdflatex
\documentclass[11pt,a4paper]{article}
\usepackage[utf8]{inputenc}
\usepackage[T1]{fontenc}
\usepackage[brazilian]{babel}
\usepackage{amsmath,amssymb,amsthm}
\usepackage[margin=2.4cm]{geometry}
\usepackage{booktabs}
\usepackage{array}
\usepackage{enumitem}
\usepackage{xcolor}
\usepackage[hidelinks]{hyperref}
\usepackage{microtype}

\sloppy
\emergencystretch=2em

\newtheorem{definicao}{Definição}
\newtheorem{teorema}{Teorema}
\newtheorem{proposicao}{Proposição}
\newtheorem{observacao}{Observação}
\newtheorem{lema}{Lema}

\begin{document}

\newcommand{\code}[1]{\texttt{\small #1}}
\newcommand{\sepcol}{\quad\penalty0$\big|$\quad\penalty0}
\newcommand{\medido}[1]{\par\medskip\noindent{\small\color{cinza}\textbf{Medido.} #1}\par\medskip}
\definecolor{cinza}{gray}{0.35}

% ─────────────────────────────────────────────────────────────────────────────
%  papers/gkcapa.tex — a capa do Reino, mínima, para os papers em `article`.
%
%  O estilo.tex completo é do `report` (títulos de capítulo, skak, …). Aqui fica
%  só o que a capa precisa, para o paper continuar a compilar sozinho com
%  pdflatex. O tradutor próprio (tests/tex) carrega o estilo.tex da raiz / o
%  symlink papers/estilo.tex e avalia o mesmo `\gkcapa`.
% ─────────────────────────────────────────────────────────────────────────────
\definecolor{ouro}{HTML}{B8912F}
\definecolor{ouroclaro}{HTML}{F3E9CE}
\definecolor{tinta}{HTML}{1E1B16}
\definecolor{regua}{HTML}{6E6858}
\providecommand{\gknota}{\fontsize{7.62}{11.05}\selectfont}
\providecommand{\gktexto}{\fontsize{10.50}{15.22}\selectfont}
\providecommand{\gksub}{\fontsize{12.33}{17.87}\selectfont}
\providecommand{\gksec}{\fontsize{14.47}{20.98}\selectfont}
\providecommand{\gkcap}{\fontsize{16.99}{24.63}\selectfont}
\providecommand{\gktit}{\fontsize{23.42}{33.95}\selectfont}
\providecommand{\Prj}{\mathbb{P}}
\providecommand{\Trf}{\mathcal{T}}
\providecommand{\gkcapa}[3]{%
\title{\vspace{-0.6cm}
{\color{ouro}\rule{\textwidth}{1.4pt}}\\[6mm]
\textbf{\gktit\color{tinta} Reino Dourado}\\[3mm]{\gksec\itshape\color{regua} Golden Kingdom}\\[8mm]
{\gkcap\scshape\color{ouro} #1}\\[5mm]
{\gksub\color{regua} #2}\\[2mm]
{\gktexto\color{regua}\itshape #3}\\[7mm]
{\color{ouro}\rule{\textwidth}{0.6pt}}\\[9mm]{\gknota\color{regua}\begin{minipage}{\textwidth}\setlength{\parindent}{0pt}%
\textbf{Aviso de Isenção Algorítmica:} O universo, as leis e as entidades contidas nestas páginas ---
incluindo felinos matriciais, aves de rapina algébricas e amuletos de controle de fluxo --- são
construções de pura ficção formal. Esta obra não descreve a realidade física, não serve como
engenharia reversa do cosmos e não constitui doutrina religiosa. Qualquer semelhança com bugs reais do seu
sistema operacional é mera coincidência (ou falta de reza). Prossiga por sua própria conta e
risco.\end{minipage}}}
\author{}\date{}}

\gkcapa{Transformada universal}
       {\textbf{Mudança reversível de leitura}}
       {círculo, reta, corda, carta local, leitura, transformada universal}
\maketitle

\begin{abstract}
O objeto aqui construído não é uma nova aritmética: é uma ponte entre duas leituras da mesma estrutura discreta. A família de polígonos regulares aproxima a circunferência em sentido geométrico apropriado; a cadeia
\[
\mathbb N\subset\mathbb Z\subset\mathbb Q\subset\mathbb R
\]
produz o contínuo linear por contagem, orientação, densidade e completamento. A ligação entre as duas é a carta local
\[
 x=R\theta,
\]
porque a distância cordal
\[
\ell(\theta_1,\theta_2)=2R\sin\left(\frac{|\theta_1-\theta_2|}{2}\right)
\]
lineariza-se localmente:
\[
\ell(\theta_1,\theta_2)=R|\theta_1-\theta_2|+O(|\theta_1-\theta_2|^3).
\]
O erro a evitar é simples: não se conclui que $P_n\to\mathbb R$. O que se conclui é mais rigoroso: a família $P_n$ realiza geometricamente o círculo; a cadeia $\mathbb Q\to\mathbb R$ produz o contínuo linear; e a carta local compatibiliza as duas escalas sem afirmar que o círculo seja a reta. A transformada universal não inventa a reta; ela organiza uma mudança reversível de leitura entre o suporte circular, a métrica cordal e a coordenada linear local. A métrica não gera probabilidade: ela lê concordância da estrutura. A distribuição só aparece depois da multiplicidade e da normalização.
\end{abstract}

\section*{Onde este documento vive}

\noindent\textbf{Camada N--R}, não a Parte Geometria nem a Parte Álgebra em sentido estrito; é a ponte entre a leitura circular e a leitura linear da mesma estrutura, sem identificar globalmente o círculo com a reta.
\[
\boxed{\;
\underbrace{\mathbb N\to\mathbb Z\to\mathbb Q\to\mathbb R}_{\text{cadeia aritmética}}
\qquad\text{e}\qquad
\underbrace{P_n\to S^1}_{\text{realização geométrica global}}\;}
\]
A carta local $x=R\theta$ não identifica os dois níveis, mas faz a passagem entre eles. Por isso a reta não é a imagem literal do círculo: é a leitura linear que surge quando a métrica cordal é observada em escala infinitesimal.

\section{Introdução}

A reta geométrica não deve ser posta como limite direto dos polígonos regulares. O motivo é simples: a família inscrita $P_n$ aproxima a circunferência, e não a reta. Em outras palavras,
\[
P_n \longrightarrow S^1
\quad\text{no sentido geométrico apropriado},\qquad S^1\neq\mathbb R.
\]
O ponto decisivo é separar suporte, leitura e métrica. O círculo e a reta não são duas fases de um mesmo objeto global; são duas leituras compatíveis da mesma estrutura em escalas diferentes. A métrica lê a concordância estrutural da representação; a distribuição só surge depois da multiplicidade e da normalização. Assim, a carta local não produz probabilidade nem converte literalmente o círculo em reta. Ela apenas lineariza a geometria local.

O que se pode fazer, de modo matematicamente correto, é distinguir duas construções distintas, no mesmo quadro da casa:

\begin{itemize}[leftmargin=1.5em]
\item a estrutura aritmética ordenada
\[
\mathbb N\subset\mathbb Z\subset\mathbb Q\subset\mathbb R,
\]
que dá contagem, orientação, densidade e completamento no registro linear;
\item a realização geométrica global
\[
P_n\longrightarrow S^1,
\]
que dá a forma circular do suporte e a sua leitura como espaço de estados.
\end{itemize}

A ponte entre essas duas construções não é uma nova ontologia, mas a carta local
\[
 x=R\theta.
\]
Com essa carta, a família de cordas produz uma métrica local que é assintoticamente linear. O que se obtém é então o seguinte:

\begin{itemize}[leftmargin=1.5em]
\item o círculo é o limite global da família de polígonos;
\item o contínuo linear aparece como leitura local e como completamento da estrutura ordenada;
\item a metricidade linear local surge da aproximação da corda por uma reta em escala pequena;
\item a transformada universal não cria uma reta global; ela compatibiliza as leituras angular e linear sem forçar uma unicidade global.
\end{itemize}

Neste ponto, o contínuo não é uma substituição imediata do discreto. Ele é a compatibilização de duas escalas: a escala global do círculo e a escala local da reta, conectadas pela carta $x=R\theta$ e pelo completamento da estrutura racional.

\section{Polígonos, circunferência e duas leituras}

Seja $P_n$ o polígono regular de $n$ lados inscrito no círculo unitário. Os vértices podem ser escritos como
\[
 z_k = e^{2\pi i k/n}, \qquad k=0,\dots,n-1.
\]
A distância entre vértices vizinhos é
\[
 d_n = |z_{k+1}-z_k| = 2\sin\left(\frac{\pi}{n}\right).
\]
Para uma corda geral que conecta $z_k$ e $z_j$, a distância é
\[
 |z_j-z_k| = 2\left|\sin\left(\frac{\pi (j-k)}{n}\right)\right|.
\]
A família de cordas contém a informação métrica do polígono. Quando $n$ cresce, as cordas se tornam pequenas e a estrutura poligonal aproxima a curva circular.

\begin{definicao}[duas leituras de uma mesma estrutura]
A mesma família de vértices $V_n=\{e^{2\pi i k/n}\}$ admite duas leituras distintas:

\begin{enumerate}
\item \textbf{leitura angular}: a ordem é cíclica e dada pelo índice $k$ módulo $n$;
\item \textbf{leitura cordal}: a métrica é dada pelos comprimentos de corda
\[
\ell_{jk}=|z_j-z_k|.
\]
\end{enumerate}

Estas duas leituras são compatíveis, mas não idênticas. O círculo carrega a leitura angular global; a reta é a realização linear local dessa leitura.
\end{definicao}

\section{A cadeia aritmética e a carta local}

A cadeia
\[
\mathbb N\subset\mathbb Z\subset\mathbb Q\subset\mathbb R
\]
fornece a leitura aritmética ordenada do contínuo, mas não é, por si só, o que transforma a família de cordas numa reta geométrica. No quadro do corpo universal, o objeto primário não é o conjunto numérico, mas o suporte que o operador descreve: dobra, lista, dualidade e realização. O que se tem, de fato, é a composição de duas construções distintas:
\[
\boxed{\underbrace{\mathbb N\to\mathbb Z\to\mathbb Q\to\mathbb R}_{\text{leitura aritmética do contínuo}}}
\qquad\text{e}\qquad
\boxed{\underbrace{P_n\to S^1}_{\text{realização geométrica global}}.}
\]
A ponte entre essas duas realizações não é uma fusão dos níveis; é a carta local
\[
\boxed{x=R\theta.}
\]
Em cada degrau da cadeia aritmética, o objeto recebe uma nova camada de compatibilização: os naturais dão a contagem, os inteiros introduzem a orientação, os racionais dão a densidade local e os reais fecham a métrica por completamento. A família de cordas, por sua vez, fornece uma realização geométrica global do mesmo sistema; a reta só aparece quando essa realização é linearizada localmente.

Para isso, tomamos uma direção orientada no círculo, isto é, um ponto base $p_0$ e um parâmetro local $\theta\in(-\varepsilon,\varepsilon)$ em torno de $p_0\in S^1$. A diferença de posições entre dois pontos vizinhos é então
\[
\Delta x = R\,\Delta \theta.
\]
A distância entre os pontos $p(\theta_1)$ e $p(\theta_2)$ é dada pela corda
\[
\ell(\theta_1,\theta_2)=2R\sin\left(\frac{|\theta_1-\theta_2|}{2}\right).
\]
Para $|\theta_1-\theta_2|\ll 1$, temos
\[
\ell(\theta_1,\theta_2)
=
R|\theta_1-\theta_2|\bigl(1+O(|\theta_1-\theta_2|^2)\bigr).
\]
Equivalentemente,
\[
\boxed{
\lim_{\Delta\theta\to0}
\frac{\ell(\theta,\theta+\Delta\theta)}{R|\Delta\theta|}=1.
}
\]
Se definimos a coordenada linear por $x(\theta)=R\theta$, então
\[
\ell(\theta_1,\theta_2)=|x(\theta_1)-x(\theta_2)|+O\!\bigl(R|\theta_1-\theta_2|^3\bigr).
\]
Este cálculo mostra que a corda é localmente assintótica à métrica linear. A construção da reta real exige, porém, uma etapa adicional: o completamento da métrica linear no registro linear escolhido pela carta. Em termos precisos,
\[
R\mathbb Q=\{Rq:q\in\mathbb Q\},\qquad \overline{R\mathbb Q}\cong\mathbb R,
\]
com $R>0$. A reta não é a imagem literal do círculo. Ela é o fechamento ordenado e métrico da leitura linear local que nasce na carta $x=R\theta$; sem essa escolha de coordenada, o círculo não se identifica com a reta.

\section{Transformada universal e mudança de leitura}

A transformada universal não é uma aplicação global
\[
S^1\longrightarrow\mathbb R.
\]
Ela também não deve ser identificada com uma única fórmula particular. O seu papel estrutural é mais geral: estabelecer uma regra reversível de mudança entre representações compatíveis do mesmo suporte. Neste quadro, a mudança de leitura não cria nenhuma nova entidade nem deriva probabilidade da métrica; ela apenas reescreve a mesma estrutura em outra coordenada. A geometria local e a distribuição estatística ficam em níveis distintos: a primeira lê concordância, a segunda nasce do peso contado e da sua normalização.

Na presente construção aparecem três leituras relacionadas:
\[
\boxed{
\text{leitura angular}
\qquad
\text{leitura linear}
\qquad
\text{leitura espectral}.
}
\]
A transformada universal é a regra que permite passar entre essas representações sem confundir o suporte com a sua representação.

Na leitura geométrica local, essa regra é realizada pela aplicação
\[
\boxed{T_R:\theta\longleftrightarrow x=R\theta,}
\]
com inversa
\[
\boxed{T_R^{-1}:x\longmapsto \theta=\frac{x}{R}.}
\]
Portanto, $T_R$ não é uma transformação global do círculo na reta. É uma mudança reversível de coordenada dentro de uma carta angular local.

A compatibilidade dessa mudança com a geometria é verificada pela corda
\[
\ell_R(\Delta\theta)=2R\left|\sin\frac{\Delta\theta}{2}\right|.
\]
Como
\[
\frac{\ell_R(\Delta\theta)}{|T_R(\Delta\theta)|}
=
\frac{2\left|\sin(\Delta\theta/2)\right|}{|\Delta\theta|}
\longrightarrow 1,
\qquad \Delta\theta\to 0,
\]
temos
\[
\boxed{\ell_R(\Delta\theta)\sim |T_R(\Delta\theta)|.}
\]
Assim, a corda não realiza a transformação: ela verifica que a mudança de leitura angular para linear preserva assintoticamente a métrica na escala local.

A mesma estrutura de mudança de leitura aparece no registro espectral. Para uma sequência $a=(a_j)$ sobre um grupo finito, define-se
\[
\boxed{
\mathcal F(a)_k
=
\frac1{\sqrt n}\sum_j a_j\chi_k(j),
}
\]
com inversa
\[
\boxed{
\mathcal F^{-1}(A)_j
=
\frac1{\sqrt n}\sum_k A_k\chi_{-k}(j).
}
\]
Aqui a transformada troca a leitura no suporte pela leitura no dual. A convolução
\[
(a*b)_j = \sum_{u+v\equiv j} a_u b_v
\]
é convertida em produto pontual,
\[
\boxed{
\mathcal F(a*b)=\sqrt n\,\mathcal F(a)\mathcal F(b).
}
\]
Quando o espectro do segundo fator não se anula, a operação inversa recupera a leitura original por
\[
\boxed{
 a = \mathcal F^{-1}\!\left(\frac{\mathcal F(c)}{\sqrt n\,\mathcal F(b)}\right).
}
\]

A partir daqui, é importante deixar explícito o nível de abstração. $T_R$ e $\mathcal F$ não são a mesma aplicação. Elas são duas realizações diferentes do mesmo princípio estrutural: a mudança reversível de leitura entre representações compatíveis do mesmo suporte. Em termos formais, a transformada universal é o princípio
\[
\boxed{
\text{representação}_1\; \overset{\mathcal T}{\longleftrightarrow}\; \text{representação}_2,
}
\]
com inversibilidade e compatibilidade preservadas. A particularização geométrica local é
\[
T_R:\theta\longleftrightarrow x=R\theta,
\]
 enquanto a particularização espectral é
\[
\mathcal F:a\longleftrightarrow A.
\]

A geometria local e o dual espectral são, portanto, duas realizações diferentes da mesma ideia estrutural: a mesma estrutura admite leituras compatíveis em coordenadas diferentes.
\[
\boxed{
\begin{array}{ccc}
\theta
&\leftrightarrow&
x
\\[6pt]
 a
&\leftrightarrow&
A
\end{array}}
\]
Na primeira linha, a mudança de leitura é entre coordenada angular e coordenada linear. Na segunda, é entre representação no suporte e representação no dual.

A corda fornece a compatibilidade métrica da primeira linha; a identidade de convolução fornece a compatibilidade algébrica da segunda. Em resumo, a arquitetura é:
\[
\boxed{
\text{corda verifica}
\quad|\quad
T_R\text{ realiza}
\quad|\quad
\mathcal F\text{ realiza no dual}.}
\]

A contribuição deste texto não é resolver a Parte Estatística em toda a sua extensão, mas fechar a dívida estrutural que a precede. O que ele resolve é a passagem
\[
\boxed{
\text{estrutura discreta}
\longrightarrow
\text{leitura}
\longrightarrow
\text{contínuo}
}
\]
sem cometer o salto ilegítimo $P_n\to\mathbb R$. Em outras palavras, ele mostra que a leitura linear não é a imagem literal do círculo, mas a realização local compatível com a métrica cordal; e que a leitura espectral não é uma nova entidade, mas outra representação da mesma estrutura. Neste quadro, a probabilidade não é tomada como princípio primeiro; ela é a normalização do peso contado sobre a estrutura. O papel da carta local é compatibilizar escalas e leituras, não produzir distribuição por si só.

\subsection{Assinatura completa e recuperabilidade do corpo}

A ideia estrutural que encerra a construção é a seguinte. Se a assinatura do corpo for um pacote fiel de invariantes, então ela fixa a classe do corpo e a sua leitura em qualquer representação compatível. Formalmente, para uma classe de corpos $\mathcal C$ e uma assinatura $\operatorname{Sig}(B)$, exigimos:
\[
\boxed{
\operatorname{Sig}(B_1)=\operatorname{Sig}(B_2)
\iff
B_1\sim B_2
}
\qquad\text{quando a assinatura é completa e fiel para a classe considerada.}
\]

Assim, a assinatura não é um rótulo arbitrário; por hipótese, ela é um pacote completo de invariantes da classe considerada, e o dado que permanece invariável sob mudança de base e sob refinamento de representação. Quando isso vale, a transformada passa a ter um papel estrutural mais profundo:
\[
\boxed{
B\overset{\mathcal F}{\longleftrightarrow}\operatorname{Sig}(B)
}
\]
com a inversa recuperando a classe do corpo somente até equivalência.

Essa condição resolve a dúvida central do texto: não se afirma que um único número contém o corpo inteiro. A assinatura precisa ser um pacote de invariantes,
\[
\operatorname{Sig}(B)=\bigl(N,\operatorname{tr},\Delta,p,q,r,\tau,\ldots\bigr),
\]
completando a informação que distingue o corpo dentro da classe considerada. Sem essa fidelidade, a inversa da transformada não produz recuperação estrutural; apenas uma leitura parcial.

A mesma ideia vale para a ponte discreto $\to$ contínuo. O limite não deve ser lido como um salto metafísico para o infinito, mas como refinamento da representação da mesma estrutura:
\[
\boxed{
B_n
\overset{\mathcal F_n}{\longleftrightarrow}
\operatorname{Sig}(B_n)
\quad\xrightarrow{\;n\to\infty\;}
\quad
B_\infty
\overset{\mathcal F_\infty}{\longleftrightarrow}
\operatorname{Sig}(B_\infty)
}
\]
com o cuidado de que o limite deve ser tomado na categoria e na estrutura adequadas. Em outras palavras, o contínuo aparece como uma leitura refinada do mesmo objeto, e não como uma entidade criada por mera passagem ao limite.

A arquitetura que resulta desta observação é então:
\[
\boxed{
\text{discreto}
\rightarrow
\text{representação}
\rightarrow
\text{assinatura}
\rightarrow
\text{refinamento}
\rightarrow
\text{contínuo}
}
\]

e em seguida,
\[
\boxed{
\text{estrutura}
\rightarrow
\begin{cases}
\text{métrica}\rightarrow\text{concordância},\\[2mm]
\text{multiplicidade}\rightarrow\text{peso}
\end{cases}
\rightarrow
\text{normalização}
\rightarrow
\text{distribuição}
\rightarrow
\text{Born}.
}
\]

O ponto essencial é que a transformada não cria o contínuo: ela preserva a estrutura enquanto a representação é refinada. A assinatura é, precisamente, o que permanece reconhecível durante essa mudança de resolução.

O mesmo princípio aparece no dual do repositório com uma forma mais precisa: no caso quadrático em que $\alpha^2=D$, a transformada universal é a avaliação nas duas folhas do corpo,
\[
\operatorname{eval}_{\pm}(a+b\alpha)=a\pm b\sqrt{D},
\]
com a multiplicação do corpo sendo convertida em convolução no espaço de coeficientes. Isso é exatamente o que o Teorema de \emph{corpo\_analitico.tex} enuncia: a ida é a convolução, a volta é a deconvolução, e a deconvolução só existe onde o espectro não se anula. O que o repositório mede em \texttt{tests/convolucao\_universal.js} confirma esse ponto: fora dos divisores de zero, a deconvolução recupera o original; quando o espectro se anula, a informação se perde e a volta falha por perda de informação, não por uma lei nova. A transformada universal, portanto, não inventa uma nova entidade: ela move a operação do suporte para o dual e a torna reversível.

A sequência que se abre a partir daqui é, portanto,
\[
\boxed{
\text{realização}
\longrightarrow
\text{representação}
\longrightarrow
\text{distribuição}
\longrightarrow
\text{Born}
}
\]
mas esse último passo não é o objeto deste paper. O que aqui fica fechado é a infraestrutura que torna legítima essa passagem: o suporte, a carta local, a reversibilidade da leitura e a compatibilidade entre geometria e dual espectral. É essa camada que torna a Estatística possível sem axioma cego de continuidade.

\section{Conclusão}

A reta geométrica não é o limite global dos polígonos. O limite global é o círculo. A reta surge como completamento da leitura linear local, habilitada pela carta $x=R\theta$ e pela compatibilização métrica da corda. A transformada universal preserva esse ponto: ela não identifica o círculo com a reta, mas organiza a passagem de uma leitura para outra sem esquecer a diferença entre o suporte circular e a leitura linear. Nesse sentido, a geometria da corda e a aritmética da reta são duas descrições compatíveis da mesma estrutura, não duas entidades confundidas em uma única imagem.

\end{document}
